\chapter[1991 --- Neher et al.]{1991: Erwin Neher, Bert Sakmann}
\label{ch:1991_Neher}

\section*{Main Contribution}

\textit{Developed the patch clamp technique to study single ion channels in cell membranes, revealing how individual protein molecules control nerve impulses.}

\section{Official Citation}

for their discoveries concerning the function of single ion channels in cells

\section{Background and Historical Context}

The ability of cells to generate and transmit electrical signals is fundamental to life. Neurons fire action potentials, muscle cells contract, and endocrine cells secrete hormones---all processes that depend on the controlled flow of ions across cell membranes. By the late twentieth century, physiologists understood that ion channels---tiny protein pores embedded in the cell membrane---were responsible for these currents. But how these channels behaved as individual molecular machines remained unknown. The development of the patch-clamp technique by Erwin Neher and Bert Sakmann in the 1970s and 1980s changed this, enabling scientists for the first time to observe the electrical activity of a single ion channel molecule in real time.

Erwin Neher was born on March 20, 1944, in Landsberg am Lech, Bavaria, Germany. He studied physics at the Technical University of Munich and then earned his doctorate at the University of Wisconsin under the supervision of Henry Lecar, working on the electrical properties of mitochondria. In 1969, he joined the laboratory of Neher at the Max Planck Institute for Psychiatry in Munich, where he first encountered the challenges of recording electrical signals from biological membranes.

Bert Sakmann was born on June 12, 1942, in Stuttgart, Germany. He studied medicine at the University of T{\"u}bingen and the University of Munich, and then completed his doctoral work at the Max Planck Institute for Psychiatry, where he and Neher first collaborated. Sakmann's medical training gave him a deep appreciation for the physiological significance of ion channels, while his experimental skills made him an ideal partner for Neher's physics-based approach.

The intellectual context of their work was shaped by the Hodgkin--Huxley model of the action potential, formulated in the early 1950s. Alan Hodgkin and Andrew Huxley had demonstrated, using the giant squid axon and the voltage-clamp technique, that the action potential arises from the sequential opening and closing of sodium and potassium channels in the nerve membrane. However, their measurements were made from large membrane patches containing thousands of channels, and the behavior of individual channel molecules remained hidden within the ensemble average. To understand how ion channels truly worked---their kinetics, selectivity, and regulation---a method capable of resolving single-channel currents was needed.

\section{The Discovery/Contribution}

\subsection{The Invention of the Patch-Clamp Technique}

Neher and Sakmann's achievement was the development of the patch-clamp technique, a method of notable technical elegance. The basic principle is deceptively simple: a glass micropipette with an extremely fine tip (less than one micrometer in diameter) is pressed against the surface of a cell, and a tight seal (``gigaseal'') is formed between the glass and the cell membrane. This seal, with a resistance of several gigaohms, electrically isolates a small patch of membrane beneath the pipette tip. When voltage is applied across the membrane patch, the current flowing through any ion channels in that patch can be measured with sufficient precision to resolve the opening and closing of individual channel molecules.

The practical realization of this technique was immensely challenging. Neher and Sakmann spent years optimizing the shape and surface properties of the glass pipettes, the composition of the bathing solution, and the electronic amplification circuitry. Early attempts yielded noisy, unstable recordings. The breakthrough came when they discovered that gentle suction applied to the pipette interior could dramatically improve the seal between glass and membrane, reducing electrical noise to a level where single-channel currents could be clearly detected.

The first successful recordings of single-channel currents were obtained from muscle cells at the neuromuscular junction. Neher and Sakmann showed that acetylcholine, the neurotransmitter released at the synapse, opened discrete ion channels in the muscle endplate membrane, each with a fixed conductance and characteristic open and closed times. These channels were not continuously open but rather flickered rapidly between open and closed states, responding stochastically to the binding and unbinding of acetylcholine molecules.

\subsection{Types of Patch-Clamp Configurations}

As the technique matured, Neher and Sakmann developed several configurations to study different aspects of channel behavior. In the cell-attached configuration, the pipette remains sealed to an intact cell, allowing channels to be studied in their native membrane environment. In the whole-cell configuration, the membrane patch under the pipette is ruptured by stronger suction, establishing electrical continuity between the pipette interior and the cell cytoplasm, enabling measurement of currents flowing through all channels in the entire cell membrane. In the outside-out and inside-out configurations, small membrane patches are excised from the cell and oriented so that the extracellular or intracellular face of the membrane faces the bathing solution, respectively. These excised patches allow precise control of the ionic environment on both sides of the membrane, enabling detailed biophysical studies of channel gating and selectivity.

\subsection{Key Discoveries}

Through their patch-clamp studies, Neher and Sakmann and their colleagues made a series of fundamental discoveries about ion channels. They showed that individual channels exhibit all-or-none behavior: a channel is either fully open or fully closed, with no intermediate states. They measured the unitary conductance of many different channel types---the amount of current that flows through a single open channel---and showed that this conductance is a fixed property of the channel protein, determined by its structure and the ions it conducts.

They demonstrated that channels are highly selective: potassium channels conduct potassium ions but not sodium or chloride ions, and vice versa. This selectivity is achieved by the channel's selectivity filter, a narrow region of the pore lined with precisely positioned amino acid residues that interact with permeating ions.

They also revealed the kinetic behavior of channels---how quickly channels open and close, how long they remain in each state, and how these properties are modulated by voltage, ligands, and intracellular signaling molecules. This kinetic information was essential for understanding the molecular mechanisms of channel gating and for building quantitative models of neuronal and cardiac electrical activity.

\section{Impact on Modern Medicine}

The patch-clamp technique revolutionized virtually every area of physiology and pharmacology. It became the gold standard for studying ion channels and has been used in thousands of laboratories worldwide. The technique's impact on medicine has been significant in several ways.

First, the patch-clamp method enabled the identification and characterization of ion channel diseases---genetic disorders caused by mutations in ion channel genes. These include cystic fibrosis (caused by mutations in the chloride channel CFTR), long QT syndrome (caused by mutations in cardiac potassium or sodium channels), certain forms of epilepsy, and various channelopathies of skeletal muscle. Understanding the biophysical defects caused by specific mutations has enabled the development of targeted therapies and diagnostic tests.

Second, the patch-clamp technique has been indispensable for drug discovery. Many pharmaceutical agents---including local anesthetics, antiarrhythmic drugs, anticonvulsants, and antihypertensives---act by blocking or modulating specific ion channels. The ability to study drug--channel interactions at the single-channel level has guided the rational design of new drugs with improved efficacy and reduced side effects.

Third, the technique has provided the foundation for our understanding of synaptic transmission and neural circuit function. By revealing how neurotransmitter-gated ion channels operate, Neher and Sakmann's work has informed research on learning, memory, and neurological and psychiatric disorders.

\section{Legacy}

Erwin Neher and Bert Sakmann shared the 1991 Nobel Prize in Physiology or Medicine for their work. Neher became director of the Department of Membrane Biophysics at the Max Planck Institute for Biophysical Chemistry in G{\"o}ttingen, where he continued to study synaptic transmission and calcium signaling. He also became a prominent advocate for science education and public engagement with science in Germany.

Sakmann became director of the Department of Cell Physiology at the same Max Planck Institute and later at the Max Planck Institute for Medical Research in Heidelberg. He continued to study synaptic channels in the central nervous system, particularly in the hippocampus and cerebral cortex, and was elected to numerous honorary positions including membership in the Leopoldina (the German National Academy of Sciences).

The patch-clamp technique remains an essential tool in neuroscience, cardiology, and pharmacology more than three decades after its development. Its influence extends to emerging fields such as optogenetics and single-molecule sequencing, where the principles of high-resolution electrical recording pioneered by Neher and Sakmann continue to inspire new technologies. Their legacy is one of technical ingenuity applied to fundamental biological questions, yielding insights that have transformed our understanding of how cells communicate and how diseases can be treated.


\section{Modern Developments}

Neher and Sakmann's patch clamp technique has enabled modern electrophysiology. Understanding of ion channels has led to drugs for epilepsy (sodium channel blockers), cardiac arrhythmias (potassium channel blockers), and pain (TRPV1 antagonists). High-throughput patch clamp screening accelerates drug discovery. Understanding of ion channel mutations has revealed the basis of channelopathies (long QT syndrome, cystic fibrosis). Optogenetics uses light-gated ion channels to control neural activity.


\section{Bibliography}

\begin{thebibliography}{9}
\bibitem{nobel-1991}
Nobel Prize Outreach, ``The Nobel Prize in Physiology or Medicine 1991,''\emph{NobelPrize.org}.\\
\url{https://www.nobelprize.org/prizes/medicine/1991/summary/}

\bibitem{lecture-1991}
Erwin Neher, ``Nobel Lecture,''\emph{NobelPrize.org}. Winner-authored primary source; downloadable from the official Nobel Prize archive.\\
\url{https://www.nobelprize.org/prizes/medicine/1991/neher/lecture/}
\end{thebibliography}
